% Generated draft: requires mathematical and source review.
\documentclass[11pt]{article}
\usepackage[margin=25mm]{geometry}
\usepackage{fontspec}
\setmainfont{Latin Modern Roman}
\usepackage{amsmath,amssymb,mathtools,unicode-math}
\setmathfont{Latin Modern Math}
\usepackage{enumitem,needspace,xurl,hyperref}
\setlength{\parindent}{0pt}
\setlength{\parskip}{5pt}
\setlength{\emergencystretch}{3em}
\newcommand{\N}{\mathbb{N}}
\newcommand{\Z}{\mathbb{Z}}
\newcommand{\Q}{\mathbb{Q}}
\newcommand{\R}{\mathbb{R}}
\newcommand{\C}{\mathbb{C}}
\newcommand{\F}{\mathbb{F}}
\newcommand{\A}{\mathbb{A}}
\newcommand{\PP}{\mathbb{P}}
\newcommand{\E}{\mathbb{E}}
\newcommand{\eps}{\varepsilon}
\newcommand{\dd}{\,\mathrm{d}}
\begin{document}
\section*{UnsolvedMath ID: 30004794}
\textbf{Source problem code:} catalogue entry
\subsection{Definitions and mathematical statement}
The supplied record attributes the following conjectural formal-degree identity to Hiraga--Ichino--Ikeda. Its mathematical setting is incomplete, so the statement below preserves the supplied target without supplying missing hypotheses.

For every essentially discrete-series representation $\pi$ of the unspecified group $G$, with extended parameter $(\phi,\rho)$, the requested identity is
\[
 d(\pi,\nu_G)
 =\frac{\dim\rho}{|S_\phi|}
 \left|\gamma(0,\phi,\operatorname{Ad}_G)\right|.
\]
Here $d(\pi,\nu_G)$ denotes the intended formal degree relative to the measure $\nu_G$; $\phi$ is the parameter and $\rho$ its enhancement in the supplied notation; $\dim\rho$ denotes the dimension of that enhancement. The notation $|S_\phi|$ denotes the order of the finite group intended by $S_\phi$, whose construction is not supplied. The term $\gamma(0,\phi,\operatorname{Ad}_G)$ denotes the intended adjoint gamma factor evaluated at zero.

The record does not specify the ground field, the class of groups $G$, the parameter correspondence, the enhancement group, the definition of $S_\phi$, or the normalization of $\nu_G$ and the gamma factor. It also does not explain the convention for formal degrees of essentially discrete-series representations, including any unitary normalization or treatment of the center. These choices are unresolved parts of the formulation. Thus the display is a schematic target requiring completion and human mathematical review, rather than a fully specified universally quantified conjecture.
\subsection{Short English statement}
The supplied target says that the formal degree of each essentially discrete-series representation with an extended parameter equals an enhancement-dimension factor times the absolute value of an adjoint gamma factor. The setting and normalizations needed to make this precise are missing. The supplied open classification is a dated curation label, not an independently verified mathematical status.
\subsection{Sources}
\begin{itemize}
\item Local UnsolvedMath registry: permanent ID $30004794$, source problem code \texttt{catalogue entry}, title ``Hiraga--Ichino--Ikeda Degrees for Essentially Discrete Series,'' in the Oberwolfach Reports Open Problems collection. The formula is attributed in the record to Hiraga--Ichino--Ikeda; no individual proposer is supplied.
\item Supplied report citation: \emph{Automorphic Forms, Geometry and Arithmetic} (2022), Oberwolfach Reports, pp.~2089--2156, DOI \texttt{10.4171/owr/2021/39}, \url{https://doi.org/10.4171/owr/2021/39}. This is also the sole supplied literature-source link, labelled ``original Oberwolfach report.'' The background describes these reports as documenting mathematical workshops and their open-problem discussions. The record separately gives proposed year 2021 and an OWR year tag of 2021; these are retained without reconciling their relation to the citation year 2022.
\item Supplied literature assessment, checked on \textbf{2026-08-22}: the record is only an introduction to the conjecture and omits a complete mathematical claim; its resolution cannot be inferred from the fragment, so it is conservatively classified as open. This is supplied dated curation derived from the OWR material, not an independently verified assessment. The background requests independent verification of time-sensitive status and cited claims.
\item The registry also labels the statement \texttt{corrected\_verified} and asserts that the standalone statement and its hypotheses were checked against the official Oberwolfach report. This is a supplied verification claim, not a verification performed for this draft; it does not resolve the omissions identified in the supplied literature assessment.
\end{itemize}
\end{document}
